\section{Capacity equality for products of balls}
\label{sec:product-orbit-saturation}
\label{subsec:product-orbit-saturation}
We extend the single-ball support-orbit obstruction to products of
balls, where each ball contributes one full row identity. Equality in
the capacity bound forces exactly one spin (Lorentz) factor for each
ball, used only by that ball's row, even when scalar ray factors are
present. The proof combines the topology of the joint support map with
the fact that each dual row depends on only one source coordinate.

\begin{theorem}[Product saturation with scalar rays]
\label{thm:product-orbit-saturation}
Let $M=\prod_{a=1}^h S^{p_a}$, where $p_a\geq1$, and $n=\sum_a p_a$.
Suppose the full row identities
\begin{equation}\label{eq:product-orbit-rows}
 1-x_a\cdot z
 =\sum_i\ip{X_i(x)}{Y_i^a(z)}
   +\sum_j\alpha_j(x)\beta_j^a(z),
 \qquad x\in M,\quad z\in S^{p_a},
\end{equation}
hold with finitely many globally labelled $C^1$ factors. The maps $X_i,Y_i^a$
take values in simple non-ray symmetric cones $K_i$, of Jordan rank $r_i$,
Peirce constant $a_i$, and real dimension $m_i$; the scalar factors are
nonnegative. Put $c_i=a_i\lfloor r_i^2/4\rfloor$. Then
\[
 \sum_i c_i\geq n.
\]
If equality holds, there are exactly $h$ non-ray factors, and after
relabelling
\[
 K_a\cong Q_{p_a+2},\qquad Y_a^d\equiv0\quad(d\ne a).
\]
Conversely, separate direct Lorentz sections, one for each ball, attain
equality. Thus the non-ray factors at equality have total dimension
$n+2h$ and total Jordan rank $2h$. If this profile fails,
then $\sum_i c_i\geq n+1$.
\end{theorem}

The first lemma determines the topology of the joint support map. As in
Section~\ref{sec:orbits}, $\mathcal I_k(V)$ is the orbit of rank-$k$
idempotents of a simple Euclidean Jordan algebra $V$.

\begin{lemma}[Product support covering]\label{lem:product-support-cover}
Assume equality in Theorem~\ref{thm:product-orbit-saturation}. Then the
positive-capacity factors have constant, balanced, complementary ranks.
Their support maps, complemented when necessary, form a finite covering
\begin{equation}\label{eq:product-orbit-cover}
 \Sigma:M\longrightarrow
       \prod_i\mathcal I_{\lfloor r_i/2\rfloor}(V_i).
\end{equation}
The topology of this covering leaves only two possible non-spin factors,
$\mathbb S_+^3$ and $\mathbb S_+^4$, whose orbits have universal covers
$S^2$ and $S^2\times S^2$. Each of these target $S^2$ factors is
assigned to a distinct $S^2$ factor of $M$: the restriction of $\Sigma$
to that source factor, lifted to the universal cover and projected onto
the target $S^2$, has degree $\pm1$, and the restriction to every other
source $S^2$ has degree zero.
\end{lemma}
\begin{proof}
At contact, put $Y_i(x)=\sum_aY_i^a(x_a)$. Termwise nonnegativity gives
$X_i(x)\circ Y_i(x)=0$. A nonnegative differentiable scalar function has
zero derivative at each zero. Hence each scalar product in
\eqref{eq:product-orbit-rows} contributes zero to the mixed derivative
at contact. Summing the rows and differentiating in the primal and dual variables
independently gives the nondegenerate product-sphere metric.
The Peirce rank bound yields
\[
 n\leq\sum_i a_i u_i v_i\leq\sum_i c_i,
 \qquad u_i=\jrank X_i(x),\quad v_i=\jrank Y_i(x).
\]
Equality forces $u_i+v_i=r_i$ and balanced ranks everywhere. Lower
semicontinuity makes each rank constant. The support-differential
calculation in Theorem~\ref{thm:global-saturation} now applies without
change: its joint differential is injective and has equal-dimensional
domain and codomain. Compactness therefore gives the finite covering
\eqref{eq:product-orbit-cover}.

Lift this covering to universal covers. The lift is a diffeomorphism.
The source universal cover is $\R^t\times\prod_{p_a\geq2}S^{p_a}$,
where $t$ counts source circles. Each spin orbit is a sphere; each
non-spin classical orbit is a balanced Grassmannian. The Albert orbit
is $\mathbb OP^2$. We exclude all non-spin possibilities except
$\mathbb S_+^3$ and $\mathbb S_+^4$.

Every positive-degree class in the mod-two cohomology of a product of
spheres has square zero. In a complex Grassmannian of order at least
three, the degree-two Schubert class instead satisfies
$\sigma_1^2=\sigma_2+\sigma_{1,1}\ne0$, omitting terms outside its
rectangle. The same relation holds with degrees doubled for quaternionic
Grassmannians; see \cite[Section~2]{SankaranSarkar2009} for the Schubert
ring and its quaternionic counterpart. These spaces are simply
connected, so the nonzero square is also present in the universal cover.
For the Albert orbit, $H^*(\mathbb OP^2;\mathbb F_2)
=\mathbb F_2[u]/(u^3)$, $|u|=8$, gives the same obstruction
\cite[Example~4.47]{Hatcher2002}.

For a balanced real Grassmannian of order $r\geq6$, both ranks $k,r-k$
are at least three. Its simply connected oriented cover has
\[
 \pi_2\bigl(\operatorname{Gr}_k^+(\R^r)\bigr)
 =\ker\bigl((\mathbb Z/2)^2\xrightarrow{+}\mathbb Z/2\bigr)
 =\mathbb Z/2.
\]
This follows from the fibration with fiber $SO(k)\times SO(r-k)$ in
$SO(r)$; the same exact sequence proves simple connectivity. Hurewicz
then gives $H_2=\mathbb Z/2$, which remains a direct summand of $H_2$ of
the universal-cover product. The source has torsion-free $H_2$.

At real order five the oriented cover is the complex quadric
$Q^3\subset\mathbb CP^4$: an oriented orthonormal pair $(u,v)$ maps
to $[u+iv]$, identifying its plane with an isotropic complex line.
Here $H^2(Q^3;\mathbb Q)$ is one-dimensional. Indeed the same homogeneous
fibration has fundamental-group map
$\mathbb Z\times\mathbb Z/2\to\mathbb Z/2$, $(b,e)\mapsto b+e$;
its kernel is $\mathbb Z$, so Hurewicz computes $H_2=\mathbb Z$.
The hyperplane class $h$ satisfies $\int_{Q^3}h^3=2$, the degree of the
quadric, and hence $h^2\ne0$. Any square-zero degree-two class on a
product containing this quadric restricts to zero on that factor.
Consequently square-zero classes cannot span the product's degree-two
rational cohomology. They do span degree two for a product of spheres,
a contradiction. This excludes real order five without a torsion
calculation in its integral cohomology ring.

The remaining non-spin possibilities are real orders three and four,
whose oriented covers are $S^2$ and
$\operatorname{Gr}_2^+(\R^4)\cong S^2\times S^2$.
These identifications, and the classical-group fibrations used above,
are compatible with the Grassmannian descriptions in
\cite[Examples~4.53--4.55]{Hatcher2002}. All remaining universal-cover
factors are therefore spheres or contractible real lines. The
indecomposable quotient of their positive-degree cohomology records
one generator in each sphere dimension. It follows that the target
sphere dimensions, with multiplicity, equal the source dimensions at
least two. Circle multiplicities agree because the finite covering
induces an inclusion of finite index on fundamental groups, whose
free-abelian ranks count circles.

Finally, in the integral cohomology of a product of simply connected
spheres, let $u_1,\ldots,u_b$ be the degree-two sphere generators. Then
\[
 \left(\sum_j b_ju_j\right)^2
   =2\sum_{j<k}b_jb_k u_ju_k.
\]
The primitive square-zero degree-two classes are exactly the $\pm u_j$.
The ring isomorphism induced by the universal-cover diffeomorphism thus
permutes these classes up to sign. Restricting to a source factor
$S^{p_a}=S^2$ gives the claimed degree assignment.
Changing the other fixed coordinates gives homotopic restrictions in the
universal cover. Hence these degrees hold, up to orientation choices, on
every \emph{$a$-sphere}, the copy of $S^{p_a}$ in $M$ on which only
$x_a$ varies.
\end{proof}

\begin{lemma}[The two real exceptions violate the row identities]
\label{lem:product-real-cylinder}
Under the hypotheses of Lemma~\ref{lem:product-support-cover}, neither
$\mathbb S_+^3$ nor $\mathbb S_+^4$ occurs.
\end{lemma}
\begin{proof}
Fix a row $d$, a point $x_d$, and a real matrix block $i$, and let
$W=\operatorname{Ran}Y_i^d(x_d)$. Row $d$ vanishes at $z=x_d$ on the whole cylinder of points
$x$ with this $d$th coordinate, so complementarity gives
\begin{equation}\label{eq:product-fixed-kernel-cylinder}
 W\subseteq\ker X_i(x)\qquad\text{whenever $x_d$ is fixed}.
\end{equation}
Suppose first that $i$ is an order-four block. Its primal support has
rank two, and the two $S^2$ rulings of its oriented Grassmannian are
assigned to distinct source factors $a,b$. If $W\ne0$ for a row
$d\ne a$, vary only $x_a$. If $\dim W=2$, the primal plane is constant.
If $\dim W=1$, the plane lies in $W^\perp$, a fixed three-space. The
inclusion
\[
 \operatorname{Gr}_2^+(W^\perp)\cong S^2
 \longrightarrow\operatorname{Gr}_2^+(\R^4)\cong S^2_+\times S^2_-
\]
has degree $\pm1$ in both projections. To see this, represent an
oriented plane by a unit simple bivector and take its self-dual and
anti-self-dual parts. For planes in a fixed three-space, the two parts
are related by an orthogonal map. Every map factoring through this
inclusion therefore has its two degrees equal up to sign. This
contradicts the degrees $(\pm1,0)$ on the assigned $a$-sphere.
Thus $Y_i^d\equiv0$ for every $d\ne a$. Applying the same argument to
the assigned $b$-sphere makes all rows zero, contradicting the
positive complementary dual rank. No order-four block remains.

For an order-three block assigned to source $a$, the primal rank is
one or two. A nonzero $W$ for $d\ne a$ either fixes its kernel line
(primal rank two), or confines its range line to
$\mathbb{RP}(W^\perp)$ of dimension at most one (primal rank one).
The lifted map on the assigned $S^2$ then has degree zero, contradicting
Lemma~\ref{lem:product-support-cover}. Hence this block is used only by row $a$.

We next show that no other non-ray block is used by row $a$. Every other
order-three block has a distinct assigned $S^2$, so the preceding
argument applies. Each remaining block is spin. Its sphere target can
be assigned to a source sphere of the same dimension with nonzero
restriction degree: in each degree at least two choose a nonzero
term in the determinant of the invertible matrix on cohomological
indecomposables of the universal covers. For circles use the invertible
rational matrix induced on fundamental groups. These assignments are
bijections; the degree-two assignments agree with the signed
permutation already proved, so none of these blocks is assigned to $a$.
Fixing row $a$ and varying such an assigned source, a nonzero dual
factor would force the spin primal support to be constant. Indeed,
a rank-two cone has only one primitive idempotent orthogonal to any
fixed nonzero complementary element. This contradicts its nonzero
degree or winding. All other non-ray factors in row $a$ vanish.

Fix all primal coordinates except $x_a$. Row $a$ is now a factorization
of the $S^2$ ball slack through one $\mathbb S_+^3$ block and finitely
many scalar rays. At each contact the ray mixed derivatives vanish,
so this block carries the full two-dimensional mixed metric. Its
projective support map is a local diffeomorphism and its ranks are
complementary one and two. Proposition~\ref{prop:no-real-three-saturation}
excludes precisely this factorization, exchanging its two arguments
if necessary. Thus no order-three block remains either.
\end{proof}

\begin{proof}[Proof of Theorem~\ref{thm:product-orbit-saturation}]
Lemma~\ref{lem:product-support-cover} gives the capacity inequality and
covering. Its remaining conclusions and
Lemma~\ref{lem:product-real-cylinder} leave only spin factors.
The universal-cover cohomology and fundamental-group counts give
exactly one target sphere for each source sphere, of the same dimension.
The determinant assignment in the preceding proof matches each target
to one source with nonzero restriction degree or winding. No splitting
of the covering as a product map is required.

Let block $i$ be assigned to source $a$, and fix a different row $d$.
If $Y_i^d(x_d)\ne0$, complementarity makes the primitive support of
$X_i(x)$ constant as $x_a$ varies, by the uniqueness of the complementary
primitive ray in a spin cone. Its assigned degree would then be zero.
Hence $Y_i^d\equiv0$ for $d\ne a$, as asserted. The direct separate
Lorentz representations give the converse. Since all capacities are
integers, failure of the equality profile gives $\sum_ic_i\geq n+1$.
\end{proof}

If the scalar ray terms vanish identically, a shorter argument shows
more: the joint support map at equality is injective. If two points
have the same primal supports, the dual rows at each point annihilate
the primal factors at the other, and the full row identities force all
their coordinates to agree. Scalar rays break this argument, which is
why the proof above keeps a covering map until it has used the row
identities on the cylinders.

For Hermitian factors over one fixed field, equality is possible only
with order-two blocks, and every source sphere then has dimension $1$,
$2$, or $4$ over $\R$, $\C$, or $\HH$, respectively. Over all symmetric cones, spin factors of
any size are allowed, but a cap $m_i\leq d$ makes equality impossible if
some source ball has $p_a+2>d$. Both statements concern globally $C^1$
full-row factorizations, not lifts whose selections exist only on
generic patches.

The Grassmannian topology and its cohomology are classical. Our
contribution is to combine them with the full row identities to classify
capacity equality and exclude sharing of a factor between rows, also
with finitely many scalar ray terms.
Aubrun, La Piana, and M\"uller-Hermes
\cite[Theorem~1]{ALM2026Lorentz} classify cones for which every positive
linear self-map factors through a finite direct sum of Lorentz cones.
Their hypothesis concerns all positive linear maps; ours concerns
one globally $C^1$ full-row slack factorization at equality in its
mixed-curvature capacity bound. Our theorem does not assert their
Lorentz-factorization property for general cones.
